% ============================================================
%   LECTURE NOTES TEMPLATE
%   Supports: theorems, proofs, code, plots, geometry, figures
% ============================================================

\documentclass[11pt, a4paper]{article}

% ── Core packages ───────────────────────────────────────────
\usepackage[T1]{fontenc}
\usepackage[utf8]{inputenc}
\usepackage{lmodern}                   % Better fonts
\usepackage[margin=2.5cm]{geometry}    % Page margins
\usepackage{microtype}                 % Better text justification

% ── Math ────────────────────────────────────────────────────
\usepackage{amsmath, amssymb, amsthm, mathtools}
\usepackage{thmtools}                  % Extended theorem tools
\usepackage{bbm}                       % \mathbbm{1} indicator
\usepackage{bm}                        % Bold math \bm{}

% ── Colors ──────────────────────────────────────────────────
\usepackage[dvipsnames, svgnames, x11names]{xcolor}

% ── Define Color Palette ────────────────────────────────────
\definecolor{TheoremBG}    {RGB}{230, 245, 230}   % soft green
\definecolor{TheoremBorder}{RGB}{ 34, 139,  34}   % forest green
\definecolor{PropBG}       {RGB}{230, 240, 255}   % soft blue
\definecolor{PropBorder}   {RGB}{ 70, 130, 180}   % steel blue
\definecolor{CorolBG}      {RGB}{255, 245, 220}   % warm cream
\definecolor{CorolBorder}  {RGB}{205, 133,  63}   % peru
\definecolor{DefBG}        {RGB}{245, 230, 255}   % lavender
\definecolor{DefBorder}    {RGB}{128,   0, 128}   % purple
\definecolor{ExBG}         {RGB}{240, 248, 255}   % alice blue
\definecolor{ExBorder}     {RGB}{ 70, 130, 180}   % steel blue
\definecolor{RemBG}        {RGB}{255, 250, 240}   % floral white
\definecolor{RemBorder}    {RGB}{169, 169, 169}   % dark gray
\definecolor{LemmaBG}      {RGB}{230, 250, 250}   % light cyan
\definecolor{LemmaBorder}  {RGB}{ 32, 178, 170}   % light sea green
\definecolor{ProblemBG}    {RGB}{230, 240, 255}   % same as Prop
\definecolor{ProblemBorder}{RGB}{ 30,  90, 200}   % royal blue
\definecolor{CodeBG}       {RGB}{ 40,  44,  52}   % dark (One Dark)
\definecolor{CodeKey}      {RGB}{198, 120, 221}   % purple keyword
\definecolor{CodeStr}      {RGB}{152, 195, 121}   % green string
\definecolor{CodeCom}      {RGB}{ 92, 100, 112}   % gray comment
\definecolor{CodeNum}      {RGB}{209, 154, 102}   % orange number
\definecolor{LinkColor}    {RGB}{ 10,  80, 160}

% ── Theorem Environments (tcolorbox) ────────────────────────
\usepackage[most]{tcolorbox}
\tcbuselibrary{theorems, skins, breakable}

%  Macro for colored theorem boxes
\newcommand{\maketheorbox}[6]{%
  % #1 = env name, #2 = display name, #3 = BG color, #4 = border color
  % #5 = title text color, #6 = counter (share with: e.g. theorem)
  \tcbset{#1style/.style={%
    enhanced, breakable,
    colback=#3, colframe=#4,
    arc=4pt,
    fonttitle=\bfseries\color{white},
    attach boxed title to top left={yshift=-\tcboxedtitleheight/2, xshift=8pt},
    boxed title style={colback=#4, colframe=#4, sharp corners},
    sharp corners=south,
    boxrule=1pt,
    top=4mm, bottom=4mm,
  }}%
}

% ── Theorem / Proposition / etc. counters ───────────────────
\newcounter{thmcounter}[section]
\renewcommand{\thethmcounter}{\thesection.\arabic{thmcounter}}

% THEOREM
\newtcbtheorem[use counter=thmcounter]{theorem}{Theorem}{
  enhanced, breakable,
  colback=TheoremBG, colframe=TheoremBorder,
  arc=4pt,
  fonttitle=\bfseries\color{white},
  attach boxed title to top left={yshift=-\tcboxedtitleheight/2, xshift=8pt},
  boxed title style={colback=TheoremBorder, colframe=TheoremBorder,
    sharp corners, fontupper=\color{white}\bfseries},
  sharp corners=south, boxrule=1pt,
  top=5mm, bottom=4mm,
}{thm}

% PROPOSITION
\newtcbtheorem[use counter=thmcounter]{proposition}{Proposition}{
  enhanced, breakable,
  colback=PropBG, colframe=PropBorder,
  arc=4pt,
  fonttitle=\bfseries\color{white},
  attach boxed title to top left={yshift=-\tcboxedtitleheight/2, xshift=8pt},
  boxed title style={colback=PropBorder, colframe=PropBorder,
    sharp corners, fontupper=\color{white}\bfseries},
  sharp corners=south, boxrule=1pt,
  top=5mm, bottom=4mm,
}{prop}

% LEMMA
\newtcbtheorem[use counter=thmcounter]{lemma}{Lemma}{
  enhanced, breakable,
  colback=LemmaBG, colframe=LemmaBorder,
  arc=4pt,
  fonttitle=\bfseries\color{white},
  attach boxed title to top left={yshift=-\tcboxedtitleheight/2, xshift=8pt},
  boxed title style={colback=LemmaBorder, colframe=LemmaBorder,
    sharp corners, fontupper=\color{white}\bfseries},
  sharp corners=south, boxrule=1pt,
  top=5mm, bottom=4mm,
}{lem}

% COROLLARY
\newtcbtheorem[use counter=thmcounter]{corollary}{Corollary}{
  enhanced, breakable,
  colback=CorolBG, colframe=CorolBorder,
  arc=4pt,
  fonttitle=\bfseries\color{white},
  attach boxed title to top left={yshift=-\tcboxedtitleheight/2, xshift=8pt},
  boxed title style={colback=CorolBorder, colframe=CorolBorder,
    sharp corners, fontupper=\color{white}\bfseries},
  sharp corners=south, boxrule=1pt,
  top=5mm, bottom=4mm,
}{cor}

% DEFINITION
\newtcbtheorem[use counter=thmcounter]{definition}{Definition}{
  enhanced, breakable,
  colback=DefBG, colframe=DefBorder,
  arc=4pt,
  fonttitle=\bfseries\color{white},
  attach boxed title to top left={yshift=-\tcboxedtitleheight/2, xshift=8pt},
  boxed title style={colback=DefBorder, colframe=DefBorder,
    sharp corners, fontupper=\color{white}\bfseries},
  sharp corners=south, boxrule=1pt,
  top=5mm, bottom=4mm,
}{def}

% EXAMPLE
\newtcbtheorem[use counter=thmcounter]{example}{Example}{
  enhanced, breakable,
  colback=ExBG, colframe=ExBorder,
  arc=4pt,
  fonttitle=\bfseries\color{white},
  attach boxed title to top left={yshift=-\tcboxedtitleheight/2, xshift=8pt},
  boxed title style={colback=ExBorder, colframe=ExBorder,
    sharp corners, fontupper=\color{white}\bfseries},
  sharp corners=south, boxrule=0.8pt,
  top=5mm, bottom=4mm,
}{ex}

% REMARK  (lighter, no box title)
\newtcbtheorem[use counter=thmcounter]{remark}{Remark}{
  enhanced, breakable,
  colback=RemBG, colframe=RemBorder,
  arc=4pt,
  fonttitle=\bfseries\color{white},
  attach boxed title to top left={yshift=-\tcboxedtitleheight/2, xshift=8pt},
  boxed title style={colback=RemBorder, colframe=RemBorder,
    sharp corners, fontupper=\color{white}\bfseries},
  sharp corners=south, boxrule=0.6pt,
  top=5mm, bottom=4mm,
}{rem}

% PROBLEM
\newtcbtheorem[use counter=thmcounter]{problem}{Problem}{
  enhanced, breakable,
  colback=ProblemBG, colframe=ProblemBorder,
  arc=4pt,
  fonttitle=\bfseries\color{white},
  attach boxed title to top left={yshift=-\tcboxedtitleheight/2, xshift=8pt},
  boxed title style={colback=ProblemBorder, colframe=ProblemBorder,
    sharp corners, fontupper=\color{white}\bfseries},
  sharp corners=south, boxrule=1pt,
  top=5mm, bottom=4mm,
}{prob}

% ── PROOF style (white background, left gray bar, italic) ───
% Matches the screenshot: "Proof of Theorem X.X" header, plain body
\tcolorboxenvironment{proof}{
  blanker, breakable,
  left=4mm,
  before skip=6pt, after skip=10pt,
  borderline west={2pt}{0pt}{gray!60},
}
% Override proof header style
\renewcommand{\proofname}{\normalfont\color{gray}\small\scshape Proof}

% ── Code listings (listings + tcolorbox) ────────────────────
\usepackage{listings}

\lstdefinestyle{CppStyle}{
  language=C++,
  backgroundcolor=\color{CodeBG},
  basicstyle=\ttfamily\small\color{white},
  keywordstyle=\color{CodeKey}\bfseries,
  stringstyle=\color{CodeStr},
  commentstyle=\color{CodeCom}\itshape,
  numberstyle=\tiny\color{CodeNum},
  numbers=left, numbersep=8pt,
  stepnumber=1,
  breaklines=true, breakatwhitespace=false,
  tabsize=4,
  showstringspaces=false,
  morekeywords={auto, nullptr, constexpr, override, final,
                uint8_t, uint16_t, uint32_t, uint64_t,
                int8_t, int16_t, int32_t, int64_t},
  frame=none,
}

\lstdefinestyle{PythonStyle}{
  language=Python,
  backgroundcolor=\color{CodeBG},
  basicstyle=\ttfamily\small\color{white},
  keywordstyle=\color{CodeKey}\bfseries,
  stringstyle=\color{CodeStr},
  commentstyle=\color{CodeCom}\itshape,
  numberstyle=\tiny\color{CodeNum},
  numbers=left, numbersep=8pt,
  stepnumber=1,
  breaklines=true, breakatwhitespace=false,
  tabsize=4,
  showstringspaces=false,
  morekeywords={self, None, True, False, as, with, yield,
                lambda, import, from, pass, assert, raise,
                except, finally, elif},
  frame=none,
}

% Wrap lstlisting in a tcolorbox for rounded corners + title
\newtcblisting{cppcode}[1][]{
  listing only, listing style=CppStyle,
  colback=CodeBG, colframe=CodeBG!80!black,
  fonttitle=\ttfamily\small\color{white},
  title=\faCode\ C++,            % remove \faCode if not using fontawesome
  sharp corners, boxrule=0.5pt,
  #1
}

\newtcblisting{pycode}[1][]{
  listing only, listing style=PythonStyle,
  colback=CodeBG, colframe=CodeBG!80!black,
  fonttitle=\ttfamily\small\color{white},
  title=Python,
  sharp corners, boxrule=0.5pt,
  #1
}

% Fallback plain code block (language-agnostic)
\newtcblisting{code}[2][]{
  listing only,
  listing options={
    style=PythonStyle,
    language=#2,
  },
  colback=CodeBG, colframe=CodeBG!80!black,
  fonttitle=\ttfamily\small\color{white},
  title=#2,
  sharp corners, boxrule=0.5pt,
  #1
}

% ── Plots and Geometry (pgfplots + tikz) ────────────────────
\usepackage{tikz}
\usepackage{pgfplots}
\pgfplotsset{compat=1.18}
\usetikzlibrary{
  arrows.meta,         % Better arrow tips
  calc,                % Coordinate calculations
  shapes.geometric,    % circles, rectangles, etc.
  positioning,         % node placement
  decorations.markings,% marks on paths
  intersections,       % path intersections
  patterns,            % hatching
  angles, quotes,      % angle labeling
  3d,                  % 3D axes
}

% Convenience: TikZ picture wrapper with caption
\usepackage{float}
\usepackage{caption}
\usepackage{subcaption}

% ── Figures / Images ────────────────────────────────────────
\usepackage{graphicx}
\usepackage{wrapfig}   % Text-wrapped figures

% Figure box style (tcolorbox)
\newtcolorbox{figurebox}[2][]{
  enhanced, breakable,
  colback=white, colframe=gray!50,
  title=#2,
  fonttitle=\bfseries\small,
  boxed title style={colback=gray!20, colframe=gray!50},
  attach boxed title to top left={yshift=-\tcboxedtitleheight/2, xshift=8pt},
  sharp corners, boxrule=0.8pt,
  #1
}

% ── Tables ──────────────────────────────────────────────────
\usepackage{booktabs}   % \toprule, \midrule, \bottomrule
\usepackage{array}
\usepackage{multirow}
\usepackage{longtable}  % Tables spanning multiple pages

% ── Hyperlinks & PDF metadata ───────────────────────────────
\usepackage{hyperref}
\hypersetup{
  colorlinks = true,
  linkcolor  = LinkColor,
  citecolor  = LinkColor,
  urlcolor   = LinkColor,
  pdftitle   = {Lecture Notes},
  pdfauthor  = {Your Name},
}

% ── Table of Contents style ─────────────────────────────────
\usepackage{titletoc}
\usepackage{titlesec}

% Section numbering depth
\setcounter{tocdepth}{2}
\setcounter{secnumdepth}{3}

% Section title styling
\titleformat{\section}
  {\large\bfseries}
  {\thesection}{1em}{}[\titlerule]

\titleformat{\subsection}
  {\normalsize\bfseries}
  {\thesubsection}{1em}{}

% ── Bibliography ────────────────────────────────────────────
\usepackage[backend=biber, style=alphabetic]{biblatex}
% \addbibresource{references.bib}   % uncomment and add your .bib

% ── Miscellaneous Useful Packages ───────────────────────────
\usepackage{enumitem}          % Customise lists
\usepackage{todonotes}         % \todo{} margin notes
\usepackage{cleveref}          % \cref{} smart cross-refs
\usepackage{siunitx}           % SI units: \SI{9.8}{\metre\per\second}
\usepackage{physics}           % \dd, \grad, \curl, \laplacian …
\usepackage{mathrsfs}          % \mathscr{F} script font

% ── Convenience Math Macros ─────────────────────────────────
\newcommand{\R}{\mathbb{R}}
\newcommand{\N}{\mathbb{N}}
\newcommand{\Z}{\mathbb{Z}}
\newcommand{\Q}{\mathbb{Q}}
\newcommand{\C}{\mathbb{C}}
\newcommand{\E}{\mathbb{E}}
\newcommand{\Prob}{\mathbb{P}}
\newcommand{\one}{\mathbbm{1}}          % indicator
\newcommand{\eps}{\varepsilon}
% \newcommand{\norm}[1]{\left\lVert#1\right\rVert}
% \newcommand{\abs}[1]{\left\lvert#1\right\rvert}
\newcommand{\inner}[2]{\left\langle#1,\,#2\right\rangle}
\newcommand{\set}[1]{\left\{#1\right\}}
\newcommand{\floor}[1]{\left\lfloor#1\right\rfloor}
\newcommand{\ceil}[1]{\left\lceil#1\right\rceil}
\DeclareMathOperator*{\argmin}{arg\,min}
\DeclareMathOperator*{\argmax}{arg\,max}
\DeclareMathOperator{\sgn}{sgn}
% \DeclareMathOperator{\tr}{tr}
% \DeclareMathOperator{\rank}{rank}
\DeclareMathOperator{\diag}{diag}
\DeclareMathOperator{\supp}{supp}

% ── Header / Footer ─────────────────────────────────────────
\usepackage{fancyhdr}
\pagestyle{fancy}
\fancyhf{}
\fancyhead[L]{\small\textit{\CourseName}}
\fancyfoot[C]{\thepage}
\renewcommand{\headrulewidth}{0.4pt}

% ── Metadata Commands (fill these per lecture) ──────────────
\newcommand{\CourseName}{Probability \& Statistics}
\newcommand{\CourseCode}{APM\_2F010\_EP}
\newcommand{\University}{Ecole Polytechnique}
\newcommand{\Lecturer}{Prof. Tardy Yoan}
\newcommand{\StudentName}{Nikolay Kharovskiy}
\newcommand{\Semester}{Fall 2026}

% ============================================================
%   DOCUMENT BEGIN
% ============================================================
\begin{document}

% ── Title Page ──────────────────────────────────────────────
\begin{titlepage}
  \centering
  \vspace*{2cm}

  {\Huge\bfseries \CourseName \par}
  \vspace{0.5cm}
  {\large \CourseCode \par}

  \vspace{1.5cm}
  \hrule height 1pt
  \vspace{0.4cm}
  {\Large\bfseries Lecture Notes \par}
  \vspace{0.4cm}
  \hrule height 1pt

  \vspace{2cm}

  \begin{tabular}{rl}
    \textbf{University:} & \University \\[4pt]
    \textbf{Lecturer:}   & \Lecturer   \\[4pt]
    \textbf{Student:}    & \StudentName \\[4pt]
    \textbf{Semester:}   & \Semester   \\[4pt]
  \end{tabular}

  \vfill

  {\small These are personal lecture notes and may contain errors.
         For the official course material consult the lecturer.}
\end{titlepage}

% ── Table of Contents ───────────────────────────────────────
\tableofcontents
\newpage

% ============================================================
%  LECTURE 1
% ============================================================

\section{Lecture 1}

\begin{theorem}{Fubini positive}{}
Let $f: X \times Y \to [0,\infty)$, $g: X \to [0,\infty)$ and $h: Y \to [0,\infty)$ be piecewise continuous. Then:
\[\int_Y \left( \int_X f(x,y) g(x) dx \right) h(y) dy = \int_X \left( \int_Y f(x,y) h(y) dy \right) g(x) dx\]

\begin{itemize}
    \item Works without integrability assumptions.
    \item Analogue: reordering double series with nonnegative terms.
\end{itemize}
\end{theorem}

\begin{theorem}{Fubini - Tonelli}{}
Let $f: X \times Y \to \mathbb{R}$ (not necessarily positive), and let $g: X \to [0,\infty)$ and $h: Y \to [0,\infty)$ be piecewise continuous. If:
\[\int_Y \left( \int_X \vert{}f(x,y)\vert{} g(x) dx \right) h(y) dy < \infty\]

Then:
\[\int_Y \left( \int_X f(x,y) g(x) dx \right) h(y) dy = \int_X \left( \int_Y f(x,y) h(y) dy \right) g(x) dx\]

\begin{itemize}
    \item Key assumption: integrable.
\end{itemize}
\end{theorem}

\subsection*{What is a probability?}

\begin{example}{Rolling a fair dice}{}
\begin{itemize}
    \item Possible outcomes $\Omega = \{1, \dots, 6\}$.
    \item All outcomes have probability $1/6$.
    \item An event is ``The result is odd''.
\end{itemize}
\end{example}

\begin{proposition}{Discrete probability space}{}
$(\Omega, \mathbb{P})$ is a discrete prob. space. $\mathbb{P}$ is called the discrete prob. on $\Omega$.
If $\forall A \subset \Omega$, $\mathbb{P}(A) = \frac{\text{Card}(A)}{\text{Card}(\Omega)}$, then $\mathbb{P}$ is called the uniform discrete probability.
\end{proposition}

\begin{example}{Non-discrete probability space}{}
Let $E = \text{``the extraction of a number in } [0, 1] \text{ uniformly at random''}$.
$E$ is modeled with the uniform probability $\mathbb{P}$ on the space $\Omega = [0,1]$.
\[\forall 0 \le a \le b \le 1, \quad \mathbb{P}([a,b]) = b - a\]

$(\Omega, \mathbb{P})$ cannot be a discrete prob. space.
\end{example}

\begin{definition}{Probability function}{}
We define a prob. function $\mathbb{P}: \mathcal{P}([0, 1]) \to [0, 1]$ s.t.
\begin{enumerate}
    \item $\mathbb{P}([0, 1]) = 1$
    \item $\mathbb{P}$ is $\sigma$-additive
    \item $\mathbb{P}([a, b]) = b - a \quad \forall a, b$
\end{enumerate}
\end{definition}

% ============================================================
%  LECTURE 2
% ============================================================

\section{Lecture 2}

\subsection{General Probability Spaces}

\begin{definition}{Sigma-algebra}{}
Let $\Omega$ be an arbitrary set. A $\sigma$-algebra $\mathscr{A}$ is a subset of $\mathcal{P}(\Omega)$ such that:
\begin{enumerate}
    \item $\emptyset \in \mathscr{A}$
    \item If $A \in \mathscr{A}$, then $A^c \in \mathscr{A}$
    \item If $(A_n)_{n \in \mathbb{N}}$ is a sequence of events in $\mathscr{A}$, then $\bigcup_{n \in \mathbb{N}} A_n$ is also in $\mathscr{A}$.
\end{enumerate}
$\sigma$-algebras are subsets of $\mathcal{P}(\Omega)$ that are closed under the basic operations on sets.
\end{definition}

\begin{example}{Sigma-algebras}{}
Let $\Omega$ be an arbitrary set and $A \subseteq \Omega$. The following are $\sigma$-algebras:
\begin{itemize}
    \item $\mathcal{P}(\Omega)$
    \item $\{\emptyset, A, A^c, \Omega\}$
\end{itemize}
\end{example}

\begin{definition}{General probability space}{}
Let $\Omega$ be an arbitrary set, $\mathscr{A} \subseteq \mathcal{P}(\Omega)$ a sigma algebra and $\mathbb{P}: \mathscr{A} \rightarrow [0,1]$. The triplet $(\Omega, \mathscr{A}, \mathbb{P})$ is called a (general) probability space if:
\begin{enumerate}
    \item $\mathbb{P}[\Omega] = 1$
    \item If $(A_n)_{n \in \mathbb{N}}$ is a sequence of disjoint events in $\mathscr{A}$, then $\mathbb{P}\left[\bigcup_{n \in \mathbb{N}} A_n\right] = \sum_{n \in \mathbb{N}} \mathbb{P}[A_n]$.
\end{enumerate}
If conditions (i) and (ii) are verified, $\mathbb{P}$ is called a (general) probability.
\end{definition}

\begin{example}{General probability space}{}
If $(\Omega, \mathbb{P})$ is a discrete probability space, then $(\Omega, \mathcal{P}(\Omega), \mathbb{P})$ is a general probability space.
\end{example}

\begin{proposition}{Basic properties of probability spaces}{}
Let $\mathbb{P}: \mathcal{P}(\Omega) \rightarrow [0,1]$ be a probability and $A, B \subseteq \Omega$. Then:
\begin{enumerate}
    \item Let $A_1, \dots, A_n$ be a finite collection of pairwise disjoint events. Then $\mathbb{P}[A_1 \cup \dots \cup A_n] = \sum_{j=1}^n \mathbb{P}[A_j]$. This property is called finite $\sigma$-additivity or additivity.
    \item If $A \subseteq B$, $\mathbb{P}[B \setminus A] = \mathbb{P}[B] - \mathbb{P}[A]$. In particular, $\mathbb{P}[A] \le \mathbb{P}[B]$ and $\mathbb{P}[A^c] = 1 - \mathbb{P}[A]$.
    \item $\mathbb{P}[A \cup B] = \mathbb{P}[A] + \mathbb{P}[B] - \mathbb{P}[A \cap B]$. In particular, $\mathbb{P}[A \cup B] \le \mathbb{P}[A] + \mathbb{P}[B]$.
\end{enumerate}
\end{proposition}

\subsection{The Borel \texorpdfstring{$\sigma$}{sigma}-algebra}

We need a $\sigma$-algebra which contains all intervals and is smaller than $\mathcal{P}(\Omega)$. The easiest choice is to take the smallest possible sigma algebra.

\begin{definition}{The Borel $\sigma$-algebra}{}
Let $I \subseteq \mathbb{R}$ be an interval. The Borel sigma algebra $\mathcal{B}(I)$ is defined as:
\[\mathcal{B}(I) = \bigcap_{\{[a,b] : a,b \in I\} \subseteq \mathscr{A} \subseteq \mathcal{P}(I)} \mathscr{A}\]

To check that it is a good definition one has to verify that $\mathcal{B}(I)$ is indeed a sigma algebra.
\end{definition}

\begin{theorem}{Uniform Probability Measure}{}
There exists a general probability space $([0, 1], \mathcal{B}([0, 1]), \mathbb{P})$ such that $\mathbb{P}[[a,b]] = b - a$ for all $a, b \in \mathbb{R}$.
\end{theorem}

\subsection{Random Variables on General Probability Spaces}

\begin{definition}{Real random variable}{}
Let $(\Omega, \mathscr{A}, \mathbb{P})$ be a probability space. A function $X: \Omega \rightarrow \mathbb{R}$ is called a real random variable if and only if for any $a, b \in \mathbb{R}$ we have that:
\[X^{-1}([a,b]) := \{\omega : X(\omega) \in [a,b]\} \in \mathscr{A}\]

A real random variable $X$ defined on a general probability space $(\Omega, \mathscr{A}, \mathbb{P})$ is discrete if $X(\Omega)$ is at most countable.
\end{definition}

\begin{proposition}{Discrete random variables as general r.v.}{}
Let $(\Omega, \mathbb{P})$ be a discrete p.s. and $X: \Omega \rightarrow \mathbb{R}$ a r.v. Then $X$ is a real random variable on $(\Omega, \mathcal{P}(\Omega), \mathbb{P})$.
\end{proposition}

\begin{remark}{Conventions}{}
\begin{itemize}
    \item If $E = \mathbb{R}$, we call $X$ a real random variable; if $E = \mathbb{R}^n$ we call $X$ a random vector.
    \item ``Let $X$ be a random variable'' has to be read as ``Let $(\Omega, \mathbb{P})$ be a discrete probability space and $X$ be a random variable defined on it''.
    \item We may denote events without brackets, and not specify the probability space. For instance: $\mathbb{P}[\{\omega : X(\omega) \in A\}] \Leftrightarrow \mathbb{P}[X \in A]$ and $\mathbb{P}[\omega] \Leftrightarrow \mathbb{P}[\{\omega\}]$.
    \item We may replace the intersection with a comma. For example: $\mathbb{P}[\{X \in A\} \cap \{Y \in B\}] \Leftrightarrow \mathbb{P}[X \in A, Y \in B]$.
\end{itemize}
\end{remark}

\subsection{The Law (Distribution) of a Random Variable}

\begin{proposition}{The law of a random variable}{}
Let $X$ be a real valued r.v. The law (distribution) $\mu_X$ of $X$ is the application $\mu_X: \mathcal{B}(\mathbb{R}) \rightarrow [0,1]$ defined by $\mu_X(C) := \mathbb{P}[X^{-1}(C)]$. In this setting, the law $(\mathbb{R}, \mathcal{B}(\mathbb{R}), \mu_X)$ is a general probability space.
\end{proposition}

\begin{remark}{Two Kinds of Random Variables}{}
In this course, a random variable $X: \Omega \rightarrow \mathbb{R}$ will be of one of two types:
\begin{itemize}
    \item \textbf{Discrete random variables:} it exists a countable set $E \subset \mathbb{R}$ and $p_X: E \rightarrow \mathbb{R}_+$ s.t. $\forall C \in \mathcal{B}(\mathbb{R})$, $\mu_X(C) = \sum_{x \in C \cap E} p_X(x)$.
    \item \textbf{Absolutely continuous random variables:} it exists $f_X: \mathbb{R} \rightarrow \mathbb{R}_+$ s.t. $\forall a < b \in \mathbb{R}$, $\mu_X([a,b]) = \int_a^b f_X(u) du$.
\end{itemize}
Fundamental intuition: $\mathbb{P}(X \in [a-\epsilon, a+\epsilon]) = \mu_X([a-\epsilon, a+\epsilon]) \approx 2\epsilon f_X(a)$. Here, $p_X$ and $f_X$ are respectively called the discrete density and the density of the random variable $X$.
\end{remark}

\subsection{Conditional Probability and Independence}

\begin{definition}{Conditional probability}{}
Let $(\Omega, \mathbb{P})$ be a discrete probability space and $A, B$ two events such that $\mathbb{P}[B] > 0$. The conditional probability of $A$ given $B$ is denoted $\mathbb{P}[A\vert{}B]$ and defined by:
\[\mathbb{P}[A\vert{}B] := \frac{\mathbb{P}[A \cap B]}{\mathbb{P}[B]}\]

\end{definition}

\begin{definition}{Independence of two events}{}
Let $(\Omega, \mathbb{P})$ be a discrete probability space and $A, B$ two events. We say that $A, B$ are independent if $\mathbb{P}[A \cap B] = \mathbb{P}[A]\mathbb{P}[B]$. If $\mathbb{P}[B] > 0$, we have the equivalent definition $\mathbb{P}[A\vert{}B] = \mathbb{P}[A]$.
\end{definition}

\begin{definition}{Joint and marginal law}{}
Let $X: \Omega \rightarrow E$, $Y: \Omega \rightarrow F$ be two r.v. The joint law $\mu_{X,Y}$ is the law of the random vector $(X, Y)$. In particular, for all $C_1, C_2 \in \mathcal{B}(\mathbb{R})$:
\[\mu_{X,Y}(C_1 \times C_2) = \mathbb{P}[(X, Y) \in C_1 \times C_2] = \mathbb{P}[X \in C_1, Y \in C_2]\]

In this context $\mu_X, \mu_Y$ are called marginal laws. We will write $p_{X,Y}$ and $f_{X,Y}$ for respectively the discrete joint density and the joint density according to the case we study.
\end{definition}

\begin{proposition}{Marginal probabilities from joint density}{}
Let $(\Omega, \mathbb{P})$ be a discrete probability space and $X: \Omega \rightarrow E$, $Y: \Omega \rightarrow F$ two random variables. Then:
\begin{align*}
    \forall y \in F, \quad p_Y(y) &= \sum_{x \in E} p_{X,Y}(x,y) \\
    \forall x \in E, \quad p_X(x) &= \sum_{y \in F} p_{X,Y}(x,y)
\end{align*}
\end{proposition}

\begin{definition}{Independence of random variables}{}
Let $(X_1, \dots, X_N)$ be real valued random variables defined all on the same p.s. $(\Omega, \mathscr{A}, \mathbb{P})$. We say that they are independent if and only if for all possible choices of $[a_1, b_1], \dots, [a_N, b_N]$ we have:
\[\mathbb{P}[X_1 \in [a_1, b_1], \dots, X_N \in [a_N, b_N]] = \prod_{i=1}^N \mathbb{P}[X_i \in [a_i, b_i]]\]

An arbitrary family $(X_i)_{i \in I}$ of random variables defined on the same probability space is independent if and only if each finite subfamily is independent.
\end{definition}

\begin{proposition}{Independence under complementation}{}
Let $(A_i)_{i \in I}$ be independent events and $I' \subseteq I$. Then the events $(B_i)_{i \in I}$ defined by
\[B_i = \begin{cases} A_i, & i \in I' \\ A_i^c, & i \in I \setminus I' \end{cases}\]

are independent.
\end{proposition}

\begin{proposition}{Independence and joint density}{}
Let $X: \Omega \rightarrow E_1, Y: \Omega \rightarrow E_2$ be two random variables defined on the same discrete probability space $(\Omega, \mathbb{P})$. Then $X$ and $Y$ are independent if and only if:
\[\forall x \in E_1, y \in E_2, \quad p_{X,Y}(x, y) = p_X(x)p_Y(y)\]

\end{proposition}

\begin{proposition}{Functions of independent random variables}{}
Let $X, Y$ be two independent random variables taking values in $E$ and $F$ respectively. Moreover, let $H, K$ be two other sets and $f: E \rightarrow H$ and $g: F \rightarrow K$ two functions. Then the random variables $f(X)$ and $g(Y)$ are independent.
\end{proposition}

\begin{proposition}{The sum of independent random variables}{}
Let $X, Y$ be two absolutely continuous and independent random variables whose densities are $f_X$ and $f_Y$. Then, the random variable $X+Y$ is absolutely continuous. Its density $f_{X+Y}$ is given by the convolution $f_X * f_Y$ of $f_X$ and $f_Y$, defined by:
\[(f_X * f_Y)(z) = \int_{-\infty}^{+\infty} f_X(x)f_Y(z-x) dx\]

It is also true in the discrete case:
\[p_{X+Y}(z) = \sum_{x \in \mathbb{R}} p_X(x)p_Y(z-x)\]

Be careful, the notation $\sum_{x \in \mathbb{R}}$ is only licit because this is the discrete case, this is highly wrong to write it in the absolutely continuous case.
\end{proposition}

% ============================================================
%  LECTURE 3
% ============================================================

\section{Lecture 3}

\subsection{Important Distributions}

\begin{definition}{Bernoulli Distribution}{}
$X$ is a Bernoulli random variable of parameter $p$ if $p_X(1) = p$ and $p_X(0) = (1 - p)$. In this case, the distribution of $X$ is the Bernoulli distribution of parameter $p$ and we write $X \sim \text{Be}(p)$.
\end{definition}

\begin{definition}{Geometric Distribution}{}
$X$ is a geometric random variable of parameter $p$ if $p_X(0) = 0$ and $p_X(n) = p(1 - p)^{n-1}$ for $n \ge 1$. In this case, the distribution of $X$ is the geometric distribution of parameter $p$ and we write $X \sim \text{Geo}(p)$.
\end{definition}

\begin{definition}{Binomial Distribution}{}
$X$ is a Binomial random variable of parameters $N$ and $p$ if
\[ p_{X}(k) = \binom{N}{k}p^{k}(1-p)^{N-k}, \quad \forall k=0,1,\dots,N. \]
In this case, the distribution of $X$ is the binomial distribution of parameters $N$ and $p$ and we write $X \sim \text{Bin}(N,p)$.
\end{definition}

\begin{definition}{Poisson Distribution}{}
$X$ is a Poisson random variable of parameter $\lambda$, if
\[ \forall n \in \mathbb{N}, \quad p_{X}(n) = \exp(-\lambda)\frac{\lambda^{n}}{n!} \]
In this case, the distribution of $X$ is the Poisson distribution of parameter $\lambda$ and we write $X \sim \text{Poi}(\lambda)$.
\end{definition}

\subsection{Relations Between Distributions}

\begin{proposition}{From Bernoulli to Binomial}{}
Let $(X_i)_{i \ge 1}$ be a sequence of i.i.d.\ random variables such that $X_1 \sim \text{Be}(p)$. Then $\sum_{i=1}^N X_i \sim \text{Bin}(N, p)$.
\end{proposition}

\begin{proof}
Let $X = \sum_{i=1}^N X_i$. We have:
\[ \mathbb{P}(X=k) = \mathbb{P}\left(\sum_{i=1}^{N} X_i = k\right) = \mathbb{P}\left(\bigcup_{K \subset \{1, \dots, N\}, |K|=k} \{X_i=1 \forall i \in K, X_i=0 \forall i \notin K\}\right) \]
By $\sigma$-additivity, since the sets are disjoint:
\[ = \sum_{|K|=k} \prod_{i \in K} \mathbb{P}(X_i=1) \prod_{i \notin K} \mathbb{P}(X_i=0) \]
\[ = \binom{N}{k} p^k (1-p)^{N-k} \]
\end{proof}

\begin{proposition}{From Bernoulli to Geometrical}{}
Let $(X_i)_{i \ge 1}$ be a sequence of i.i.d.\ random variables such that $X_1 \sim \text{Be}(p)$. Then $T = \inf\{i \ge 1 : X_i = 1\}$ is a geometric random variable of parameter $p$, i.e., $T \sim \text{Geo}(p)$.
\end{proposition}

\begin{proof}
First, we find the probability that $T \ge n$:
\[ \mathbb{P}(T \ge n) = \mathbb{P}(X_i = 0 \; \forall i \in \{1, \dots, n-1\}) = (1-p)^{n-1} \]
Then, using the difference of probabilities:
\[ \mathbb{P}(T = n) = \mathbb{P}(T \ge n) - \mathbb{P}(T \ge n+1) = (1-p)^{n-1} - (1-p)^n \]
\[ = (1-p)^{n-1} (1 - (1-p)) = p(1-p)^{n-1} \]
\end{proof}

\begin{proposition}{Law of Small Numbers (simplified version)}{}
Let $\lambda > 0$ and $(X_n)_{n \ge \lceil\lambda\rceil}$ be a sequence of random variables such that $X_n \sim \text{Bin}(n, \lambda/n)$. In particular, this means $\mathbb{E}[X_n] = \lambda$ for all $n$. Then we have
\[ \lim_{n \to +\infty} p_{X_n}(k) = e^{-\lambda} \frac{\lambda^k}{k!} = p_X(k), \text{ with } X \sim \text{Poi}(\lambda). \]
\end{proposition}

\begin{proof}
Let $p = \lambda/n$. For a given $r$:
\[ \mathbb{P}(X_n = r) = \binom{n}{r} \left(\frac{\lambda}{n}\right)^r \left(1-\frac{\lambda}{n}\right)^{n-r} \]
\[ = \frac{n(n-1)\dots(n-r+1)}{r!} \frac{\lambda^r}{n^r} \left(1-\frac{\lambda}{n}\right)^n \left(1-\frac{\lambda}{n}\right)^{-r} \]
As $n \to \infty$, we analyze the terms:
\begin{itemize}
    \item $\frac{n(n-1)\dots(n-r+1)}{n^r} \to 1$
    \item $\left(1-\frac{\lambda}{n}\right)^n \to e^{-\lambda}$
    \item $\left(1-\frac{\lambda}{n}\right)^{-r} \to 1$
\end{itemize}
Thus, $\mathbb{P}(X_n = r) \to e^{-\lambda} \frac{\lambda^r}{r!}$.
\end{proof}

\subsection{Expectation}

\begin{definition}{Expectation}{}
Let $X$ be a non negative valued random variable defined on a general probability space $(\Omega, \mathbb{P})$. We write $\mathbb{E}[X]$ the expectation of $X$ which equals:
\begin{itemize}
    \item $\sum_{x \ge 0} x p_X(x)$ in the discrete case,
    \item $\int_{\mathbb{R}} x f_X(x) dx$ in the absolute continuous case.
\end{itemize}
If $X$ is a real valued random variable, we say that the expectation $\mathbb{E}[X]$ of $X$ exists if one of the following conditions $\mathbb{E}[\mathbbm{1}_{X \ge 0} X] < +\infty$ or $\mathbb{E}[\mathbbm{1}_{X \le 0} (-X)] < +\infty$ is verified. In that case the expectation value of $X$ is given by
\[ \mathbb{E}[X] = \mathbb{E}[\mathbbm{1}_{X \ge 0} X] - \mathbb{E}[\mathbbm{1}_{X \le 0} (-X)]. \]
\end{definition}

\begin{example}{Expectation Examples}{}
\begin{itemize}
    \item \textbf{Fair dice:} If $p_X(x) = 1/6$ for $x \in \{1, 2, 3, 4, 5, 6\}$, then $\mathbb{E}[X] = 3.5$.
    \item \textbf{Uniform law on $[0, 1]$:} If $f_X = \mathbbm{1}_{[0,1]}$, then $\mathbb{E}[X] = 0.5$.
\end{itemize}
\end{example}

\begin{proposition}{Expectation - Equivalent Definition}{}
Let $X: \Omega \to E$ be a discrete random variable defined over a discrete probability space $(\Omega, \mathbb{P})$. We have that $\mathbb{E}[X]$ exists if and only if one of the two conditions
\[ \sum_{\omega \in \Omega, X(\omega) \ge 0} X(\omega) \mathbb{P}[\{\omega\}] < +\infty, \quad \sum_{\omega \in \Omega, X(\omega) < 0} X(\omega) \mathbb{P}[\{\omega\}] > -\infty \]
is verified. In this case
\[ \mathbb{E}[X] = \sum_{\omega \in \Omega} X(\omega) \mathbb{P}[\{\omega\}] \]
\end{proposition}

\begin{proof}
Assume $X \ge 0$. Using the positive Fubini theorem:
\[ \sum_{\omega \in \Omega} X(\omega) \mathbb{P}[\{\omega\}] = \sum_{x \in X(\Omega)} \sum_{\omega \in \Omega, X(\omega)=x} X(\omega) \mathbb{P}[\{\omega\}] \]
\[ = \sum_{x \in X(\Omega)} x \sum_{\omega: X(\omega)=x} \mathbb{P}[\{\omega\}] = \sum_{x \in X(\Omega)} x \mathbb{P}(X=x) = \mathbb{E}[X] \]
The same logic extends to $\mathbb{E}[\mathbbm{1}_{X \ge 0}X]$ and $\mathbb{E}[\mathbbm{1}_{X \le 0}(-X)]$.
\end{proof}

\begin{proposition}{Expectation of a Function of a Random Variable}{}
Let $X: \Omega \to E$ be a random variable defined over a general probability space $(\Omega, \mathbb{P})$ and $g: E \to \mathbb{R}$ any function. 
In the discrete case, $\mathbb{E}[g(X)]$ exists if and only if either $\sum_{x: g(x) \ge 0} g(x) p_X(x) < +\infty$ or $\sum_{x: g(x) < 0} g(x) p_X(x) > -\infty$ is verified. Then:
\[ \mathbb{E}[g(X)] = \sum_{x \in X(\Omega)} g(x) p_X(x) \]
In the absolutely continuous case, if $g$ is piecewise continuous, $\mathbb{E}[g(X)]$ exists if and only if either $\int_{\mathbb{R}} \mathbbm{1}_{g(x) \ge 0} g(x) f_X(x) dx < +\infty$ or $\int_{\mathbb{R}} \mathbbm{1}_{g(x) \le 0} g(x) f_X(x) dx > -\infty$ is verified. Then:
\[ \mathbb{E}[g(X)] = \int_{\mathbb{R}} g(x) f_X(x) dx \]
\end{proposition}

\begin{proof}
For the discrete case, assuming $g \ge 0$:
\[ \sum_{x \in X(\Omega)} g(x) p_X(x) = \sum_{q \ge 0} \sum_{x: g(x)=q} q \mathbb{P}(X=x) \]
\[ = \sum_{q \ge 0} q \mathbb{P}(g(X)=q) = \mathbb{E}[g(X)] \]
\end{proof}

\begin{proposition}{Product of Independent Random Variables}{}
Let $X, Y$ be two independent random variables. If $\mathbb{E}[X], \mathbb{E}[Y] < +\infty$, then $\mathbb{E}[XY] < +\infty$ as well and
\[ \mathbb{E}[XY] = \mathbb{E}[X]\mathbb{E}[Y] \]
\end{proposition}

\begin{proof}
By Fubini positive, in the discrete case:
\[ \mathbb{E}[XY] = \sum_{x} \sum_{y} xy \mathbb{P}(X=x, Y=y) = \sum_{x} \sum_{y} xy \mathbb{P}(X=x)\mathbb{P}(Y=y) \]
\[ = \left( \sum_{x} x \mathbb{P}(X=x) \right) \left( \sum_{y} y \mathbb{P}(Y=y) \right) = \mathbb{E}[X]\mathbb{E}[Y] \]
\end{proof}

\begin{proposition}{Properties of Expectation}{}
Let $X, Y$ two real valued random variables defined on the same general probability space $(\Omega, \mathbb{P})$. The following holds:
\begin{enumerate}
    \item \textbf{Monotonicity:} If both $\mathbb{E}[X], \mathbb{E}[Y]$ exist, and $X(\omega) \le Y(\omega) \; \forall \omega \in \Omega$, then $\mathbb{E}[X] \le \mathbb{E}[Y]$.
    \item \textbf{Absolute Value:} If $\mathbb{E}[X]$ exists, then $|\mathbb{E}[X]| \le \mathbb{E}[|X|]$.
    \item \textbf{Linearity:} If $X, Y$ have finite expectation, and $a, b \in \mathbb{R}$, then $aX + bY$ has finite expectation and $\mathbb{E}[aX + bY] = a\mathbb{E}[X] + b\mathbb{E}[Y]$.
\end{enumerate}
\end{proposition}

\begin{proof}
For discrete random variables:
\begin{enumerate}
    \item $\mathbb{E}[X] = \sum_{\omega} X(\omega) \mathbb{P}[\{\omega\}] \le \sum_{\omega} Y(\omega) \mathbb{P}[\{\omega\}] = \mathbb{E}[Y]$.
    \item $|\mathbb{E}[X]| = \left| \sum_{\omega} X(\omega) \mathbb{P}[\{\omega\}] \right| \le \sum_{\omega} |X(\omega)| \mathbb{P}[\{\omega\}] = \mathbb{E}[|X|]$ (by triangle inequality).
    \item $\mathbb{E}[aX + bY] = \sum_{\omega} (aX(\omega) + bY(\omega)) \mathbb{P}[\{\omega\}] = a\mathbb{E}[X] + b\mathbb{E}[Y]$.
\end{enumerate}
\end{proof}

\begin{example}{Uniform Random Variables Moments}{}
If $X$ is a uniform random variable over the interval $[a, b]$, then $f_X(x) = \frac{1}{b-a}$ for $x \in [a,b]$ and $0$ otherwise.
\[ \mathbb{E}[X] = \frac{1}{b - a} \int_a^b x dx = \frac{b^2 - a^2}{2(b - a)} = \frac{a + b}{2} \]
\[ \mathbb{E}[X^2] = \frac{1}{b - a} \int_a^b x^2 dx = \frac{b^3 - a^3}{3(b - a)} = \frac{b^2 + ab + a^2}{3} \]
\end{example}

\subsection{Variance and Covariance}

\begin{definition}{Variance}{}
Let $X$ be a real valued random variable with $\mathbb{E}[X^2] < +\infty$. The variance of $X$ is
\[ \mathbb{V}(X) = \mathbb{E}[(X - \mathbb{E}[X])^2] \]
\end{definition}

\begin{definition}{Covariance}{}
Let $X, Y$ be two random variables s.t. $\mathbb{E}[|X|], \mathbb{E}[|Y|] < +\infty$ and $\mathbb{E}[|XY|] < +\infty$. The covariance of $X$ and $Y$ is
\[ \text{Cov}(X, Y) = \mathbb{E}[(X - \mathbb{E}[X])(Y - \mathbb{E}[Y])] \]
In particular, $\mathbb{V}(X) = \text{Cov}(X, X)$.
\end{definition}

\begin{proposition}{Properties of Variance and Covariance}{}
Let $X, Y$ be two random variables such that $\mathbb{E}[X^2], \mathbb{E}[Y^2] < +\infty$. Then we have
\[ \mathbb{V}(X) = \mathbb{E}[X^2] - \mathbb{E}[X]^2, \quad \text{Cov}(X, Y) = \mathbb{E}[XY] - \mathbb{E}[X]\mathbb{E}[Y]. \]
and
\[ \mathbb{V}(X + Y) = \mathbb{V}(X) + \mathbb{V}(Y) + 2\text{Cov}(X, Y). \]
\begin{itemize}
    \item If $X \sim \text{Be}(p)$, then $\mathbb{V}(X) = p - p^2$.
    \item If $X$ is uniformly distributed on $\{-N, -N + 1, \dots, N - 1, N\}$, then $\mathbb{V}(X) = \frac{N(N + 1)}{3}$.
\end{itemize}
\end{proposition}

% \subsection{\texorpdfstring{$L^p$}{Lp} Spaces}

% \begin{definition}{$L^p$ spaces and norms}{}
% For any $p \ge 0$ we set
% \[ L^p(\Omega, \mathbb{P}) = \{X : \Omega \to \mathbb{R} : \mathbb{E}[|X|^p] < +\infty\} \]
% Whenever there is no ambiguity, we abbreviate $L^p(\Omega, \mathbb{P})$ with $L^p$. If $X \in L^p$, its $p$-norm $\|X\|_p$ is defined as
% \[ \|X\|_p = (\mathbb{E}[|X|^p])^{1/p} \]
% \end{definition}

% \begin{theorem}{Seminorm property}{}
% $\| \cdot \|_p$ is a seminorm on the space $L^p(\Omega, \mathbb{P})$. Recall that it verifies homogeneity ($\|\lambda x\| = |\lambda|\|x\|$) and the triangular inequality ($\|x + y\| \le \|x\| + \|y\|$).
% \end{theorem}

% \begin{proposition}{Properties of $L^p$ spaces}{}
% Let $(\Omega, \mathbb{P})$ be a general probability space and $1 \le p \le q < +\infty$. Then $L^p(\Omega, \mathbb{P})$ is a vector space:
% \begin{itemize}
%     \item $X \in L^p \implies \lambda X \in L^p$
%     \item $X, Y \in L^p \implies X + Y \in L^p$
% \end{itemize}
% If $(\Omega, \mathbb{P})$ is a discrete probability space:
% \begin{enumerate}
%     \item If $0 < p \le q < +\infty$, $L^p(\Omega, \mathbb{P}) \supseteq L^q(\Omega, \mathbb{P})$.
%     \item If $X, Y \in L^2(\Omega, \mathbb{P})$, then $XY \in L^1(\Omega, \mathbb{P})$.
% \end{enumerate}
% \end{proposition}

\end{document}